\documentclass[UTF8,fontset=fandol,11pt]{ctexart}
\usepackage[a4paper,margin=25mm]{geometry}
\usepackage{enumitem}
\usepackage{hyperref}
\hypersetup{pdfauthor={},pdftitle={AI工具使用详情}}
\title{AI工具使用详情}
\author{}
\date{}
\begin{document}
\maketitle
\section{记录范围与工具}
本文件记录截至2026年9月13日、当前可见对话和项目文件所能确认的AI辅助使用情况。工具为Codex桌面应用中的AI助手；当前会话系统标识为GPT-6，历史各次调用的具体模型版本未完整记录，不能据此推定全部历史调用使用相同版本。本记录不预填后续尚未发生的使用行为。

\section{用途与主要提示过程}
\begin{enumerate}[leftmargin=*]
\item 模板与论文组织。用户要求部署并保持现有论文模板，整理章节结构、符号和单位，调整图表引用、程序展示与编译。助手参与相应文本和文件操作。
\item 四问正文撰写。用户逐问提供模型、数值结果和审计反馈，要求依据这些材料起草、审查和修改问题重述、问题分析、模型假设及各问正文。用户明确要求统一摄氏温度与Arrhenius绝对温度记法，区分建模假设、数值实现和题目条件，不允许编造结果。
\item 模型口径核对。用户确认固定半径问题与径向收缩问题的区别，冻结第三、四问全域临界时间和规则采样口径，并限定正式验证证据的范围。助手依据确认内容修改文字，不将数值一致性写成实验验证。
\item 结构与篇幅调整。用户要求四问采用“具体模型的建立—模型的求解—结果的得出与分析”的三段式结构，精简重复文字和非核心符号。摘要及附录曾按用户要求暂缓；本次审查阶段暂缓摘要，后续在用户授权下完成题目、摘要和关键词撰写。
\item 图件与程序附录。助手参与绘图脚本审查、路径问题识别、中文文件头注释和附录收录。本次按用户授权用既有温度、水分结果及附件1补全问题一环境响应曲线，没有重新求解主模型。
\item 本轮审查提示。用户提供四问与附录版审查意见，要求采用经讨论认可的第一、第二组建议并逐项报告修改；涉及数据范围、模型局限、表格位置、精度表述、引图删重和AI材料位置等。
\end{enumerate}

\section{采纳、修改与核验}
AI文本经用户多轮反馈修订。采纳内容包括统一术语和温标、四问模型递进、固定位置与移动表面的区分、临界时刻与采样时刻的区分，以及图文一致性修正。用户要求删除的问题一独立交叉验证叙述、问题三低湿分支诊断及相关图件不作为本版正文验证证据。

本轮还根据四问正式结果完成题目、摘要和关键词，核对了附件1的时间节点、最终支撑压缩包内第四问结果表头、正文引用和模板编译情况。正文主要数值应追溯至原始附件、结果工作簿及正式验证输出；新增的浓度因子解释仅由已有经验公式直接计算。AI辅助检查和程序语法检查均不等同于独立实验验证，也不保证全部支撑程序已在新目录下完成独立运行。

\section{记录边界}
其他协作者是否使用AI、使用何种工具及其提示过程，当前可见材料不足以完整确认，因此不作推断。本文件不是原始对话逐字导出；如后续增加AI使用、取得历史版本记录或新增协作者使用详情，应据实际情况补充。正式提交前仍需由论文负责人核对记录完整性及实际提交文件。
\end{document}
